\documentclass[../../../../uem_mei_q.tex]{subfiles}

\begin{document}
    $x$についての関数$f(x)$を\\
    \hspace{3\zw}%
    $\displaystyle%
        f(x) = x^3-3x^2+4x+3%
    $\\
    により定める．座標平面上で，曲線 $y=f(x)$ の接線で点$\mathrm{A}(1,\;3)$を通るようなものの%
    方程式を求めたい．接点を$\mathrm{T}(t,\;f(t))$とするとき，Tにおける $y=f(x)$ の接線の方程式%
    を$t$を用いて表すと\\
    \hspace{3\zw}%
    $\displaystyle%
        \begin{aligned}
        y=&\left(\:\text{\myboxb{　オ　}}\:t^2-\:\text{\myboxb{　カ　}}\:t%
        +\:\text{\myboxb{　キ　}}\:\right)x \\
        &+\left(-\:\text{\myboxb{　ク　}}\:t^3+\:\text{\myboxb{　ケ　}}\:t^2+%
        \:\text{\myboxb{　コ　}}\:\right)
        \end{aligned}
    $\\
    であり，この方程式が点Aを通ることから $t=\text{\myboxb{　サ　}}\:$ であって，接線の方程式は\\
    \hspace{3\zw}%
    $\displaystyle%
            y=\:\text{\myboxb{　シ　}}\:x-\:\text{\myboxb{　ス　}}%
    $\\
    である．
\end{document}
